%------------------------------------------------------------------------------
% \Conclusion
%------------------------------------------------------------------------------
\section*{\vspace*{-\baselineskip}\begin{center}ЗАКЛЮЧЕНИЕ\end{center}\vspace*{-\baselineskip*2}}\par
\addcontentsline{toc}{section}{\hspace*{-1em}Заключение}\par

%------------------------------------------------------------------------------
% Text
%------------------------------------------------------------------------------
\hspace*{12.5 mm}В рамках данной дипломной работы была спроектирована, 
разработана и протестирована аппаратная реализация нейронной сети на базе 
программируемой логической интегральной схемы ZYBO Z7 с использованием 
платформы PYNQ.

На первом этапе работы была произведена формализация структуры нейронной сети в 
виде управляющего и операционного автоматов, что позволило эффективно 
спроектировать архитектуру IP-блока. Для обеспечения взаимодействия между 
процессорной системой и аппаратной логикой применялась среда PYNQ, 
обеспечивающая удобную интеграцию аппаратных модулей и программного управления.

Был разработан специализированный IP-блок, реализующий основные функции 
нейронной сети, включая обработку входных данных, вычисление выходов нейронов и
применение функций активации. Особое внимание было уделено оптимизации 
аппаратной реализации: использовалась фиксированная точность представления 
данных для уменьшения занимаемых ресурсов без существенного ухудшения качества 
классификации.

Разработанная система прошла всестороннее тестирование, включающее: 

– тестирование на стандартном наборе данных MNIST с построением матрицы 
спутывания и сравнением с эталонной моделью. 

– тестирование на изображениях, созданных вручную, с целью оценки 
устойчивости работы. 

– исследование влияния разрядности весовых коэффициентов на точность 
классификации. 

Результаты экспериментов показали высокую точность классификации, сравнимую 
с полносвязными архитектурами нейронной сети, при значительно меньших затратах 
вычислительных ресурсов благодаря LST преобразованию. 

Таким образом, поставленные цели дипломной работы были достигнуты, а все задачи 
успешно решены. Разработанный IP-блок нейронной сети может быть использован в 
составе встроенных систем реального времени, требующих высокой скорости 
обработки данных при ограниченных аппаратных ресурсах.

   